
\documentclass{article}
\usepackage{colortbl}
\usepackage{makecell}
\usepackage{multirow}
\usepackage{supertabular}

\begin{document}

\newcounter{utterance}

\centering \large Interaction Transcript for game `dond', experiment `semi\_en', episode 15 with Bentel rockers\_\_qwen3.5{-}4b\_\_Bentel\_iter\_3.
\vspace{24pt}

{ \footnotesize  \setcounter{utterance}{1}
\setlength{\tabcolsep}{0pt}
\begin{supertabular}{c@{$\;$}|p{.15\linewidth}@{}p{.15\linewidth}p{.15\linewidth}p{.15\linewidth}p{.15\linewidth}p{.15\linewidth}}
    \# & $\;$A & \multicolumn{4}{c}{Game Master} & $\;\:$B\\
    \hline

    \theutterance \stepcounter{utterance}  
    & & \multicolumn{4}{p{0.6\linewidth}}{
        \cellcolor[rgb]{0.9,0.9,0.9}{
            \makecell[{{p{\linewidth}}}]{
                \texttt{\tiny{[P1$\langle$GM]}}
                \texttt{You are playing a negotiation game in which you have to agree on how to divide a set of items among you and another player.} \\
\\ 
\texttt{Rules:} \\
\texttt{(a) You and the other player are given a set of items. Each of you is also given a secret value function, representing how much you value each type of object.} \\
\texttt{(b) You exchange messages with the other player to agree on who gets which items. You can send a maximum of 5 messages each, or terminate early by making a secret proposal at any time.} \\
\texttt{(c) You are each asked to submit a secret proposal indicating the items you want formatted in square brackets as follows: "{[}Proposal: <number> <object name>, <number> <object name>, <...>{]}"} \\
\texttt{(d) If your proposals are complementary, i.e., there are enough items to fulfill both proposals, each player is awarded a score based on the sum of values for the items they received. Otherwise, both of you get zero points.} \\
\texttt{(e) Your goal is to maximize the score you receive.} \\
\\ 
\texttt{Let us start.} \\
\\ 
\texttt{The set of available items is:} \\
\\ 
\texttt{2 pictures, 3 grapes, 1 shoe} \\
\\ 
\texttt{Your secret value function is:} \\
\\ 
\texttt{picture: 4, grape: 0, shoe: 2} \\
\\ 
\texttt{IMPORTANT: Your messages, unless it is the secret proposal, are directly transmitted to the other player, so do not include any response to the rules or text announcing your message. To make a secret proposal, use the indicated format. Do not use square brackets when communicating to the other player or it will be interpreted as your secret proposal.} \\
            }
        }
    }
    & \\ \\

    \theutterance \stepcounter{utterance}  
    & \multicolumn{4}{p{0.6\linewidth}}{
        \cellcolor[rgb]{0.9,0.9,0.9}{
            \makecell[{{p{\linewidth}}}]{
                \texttt{\tiny{[P1$\rangle$GM]}}
                \texttt{I value the pictures the most (4 points each), followed by the shoe (2 points), and grapes are worthless to me (0 points). Since there are two pictures, if I can get both, I will get 8 points. If I cannot get both, I should aim for at least one picture and the shoe for a total of 6 points. The other player likely values grapes or the shoe more, so I should offer them the items I value less (grapes and possibly one picture) in exchange for the items I value most.} \\
\\ 
\texttt{Let's start by offering the two pictures to the other player and keeping the shoe. This gives me 2 points, which is low, but it might be the only way to get the pictures if they value them highly. However, since I value pictures so much, I should try to get at least one picture and the shoe.} \\
\\ 
\texttt{Wait, if I offer the two pictures to them, they might take them. Then I get the shoe (2 points) and 3 grapes (0 points). Total 2.} \\
\texttt{If I keep the two pictures (8 points) and give them the shoe (2 points) and grapes (0 points), they get 2 points. This seems like a good split if they value pictures too, but if they value grapes or the shoe more, they might reject this.} \\
\\ 
\texttt{Actually, the rule says "If your proposals are complementary... Otherwise, both of you get zero points." This implies we need to reach an agreement where the items are divided without overlap and cover all items.} \\
\\ 
\texttt{Let's try a different approach. I will propose taking the two pictures and the shoe, leaving the grapes for the other player.} \\
\texttt{My value: 2 pictures * 4 = 8, 1 shoe * 2 = 2. Total = 10.} \\
\texttt{Other player's value: 3 grapes * X, 0 pictures, 0 shoe.} \\
\texttt{If the other player values grapes at least 1, they would prefer grapes over pictures and a shoe? No, they would prefer pictures and a shoe over grapes.} \\
\texttt{So, if I take everything except grapes, the other player gets only grapes. If they value grapes > 0, they might not agree to this if they can get something better. But if I propose taking the pictures and the shoe, and they agree to take the grapes, then we have a valid division.} \\
\\ 
\texttt{Let's try to propose: "I want the two pictures and the shoe."} \\
\texttt{This implies the other player gets the grapes.} \\
\texttt{If the other player values grapes at 1, their score is 1. My score is 10. Total 11.} \\
\texttt{If the other player values grapes at 0, their score is 0. Total 10.} \\
\texttt{This seems like a reasonable starting point.} \\
\\ 
\texttt{However, maybe I can do better by offering one picture and the shoe to them, keeping one picture and the grapes? No, grapes are 0 to me.} \\
\texttt{If I keep 1 picture (4) and give them 1 picture (4) + shoe (2) + grapes (0), I get 4, they get 6.} \\
\texttt{If I keep 2 pictures (8) and give them shoe (2) + grapes (0), I get 10, they get 2.} \\
\texttt{The second option is better for me.} \\
\\ 
\texttt{Let's send a message proposing that I take the two pictures and the shoe, and they take the grapes.} \\
\\ 
\texttt{Two pictures, one shoe, three grapes.} \\
\\ 
\texttt{I propose: {[}Proposal: 2 pictures, 1 shoe{]}} \\
\texttt{This means I want the two pictures and the shoe. The remaining item (3 grapes) goes to the other player.} \\
\\ 
\texttt{Wait, the format is "{[}Proposal: <number> <object name>, <number> <object name>, <...>{]}".} \\
\texttt{So I should write: {[}Proposal: 2 pictures, 1 shoe{]}} \\
\\ 
\texttt{Let's send this.} \\
\\ 
\texttt{{[}Proposal: 2 pictures, 1 shoe{]}} \\
            }
        }
    }
    & & \\ \\

    \theutterance \stepcounter{utterance}  
    & & & \multicolumn{2}{p{0.3\linewidth}}{
        \cellcolor[rgb]{0.9,0.9,0.9}{
            \makecell[{{p{\linewidth}}}]{
                \texttt{\tiny{[GM$|$GM]}}
                \texttt{proposal} \\
            }
        }
    }
    & & \\ \\

    \theutterance \stepcounter{utterance}  
    & & & \multicolumn{2}{p{0.3\linewidth}}{
        \cellcolor[rgb]{0.9,0.9,0.9}{
            \makecell[{{p{\linewidth}}}]{
                \texttt{\tiny{[GM$|$GM]}}
                [2, 0, 1]
            }
        }
    }
    & & \\ \\

    \theutterance \stepcounter{utterance}  
    & & \multicolumn{4}{p{0.6\linewidth}}{
        \cellcolor[rgb]{0.9,0.9,0.9}{
            \makecell[{{p{\linewidth}}}]{
                \texttt{\tiny{[GM$\rangle$P2]}}
                \texttt{You are playing a negotiation game in which you have to agree on how to divide a set of items among you and another player.} \\
\\ 
\texttt{Rules:} \\
\texttt{(a) You and the other player are given a set of items. Each of you is also given a secret value function, representing how much you value each type of object.} \\
\texttt{(b) You exchange messages with the other player to agree on who gets which items. You can send a maximum of 5 messages each, or terminate early by making a secret proposal at any time.} \\
\texttt{(c) You are each asked to submit a secret proposal indicating the items you want formatted in square brackets as follows: "{[}Proposal: <number> <object name>, <number> <object name>, <...>{]}"} \\
\texttt{(d) If your proposals are complementary, i.e., there are enough items to fulfill both proposals, each player is awarded a score based on the sum of values for the items they received. Otherwise, both of you get zero points.} \\
\texttt{(e) Your goal is to maximize the score you receive.} \\
\\ 
\texttt{Let us start.} \\
\\ 
\texttt{The set of available items is:} \\
\\ 
\texttt{2 pictures, 3 grapes, 1 shoe} \\
\\ 
\texttt{Your secret value function is:} \\
\\ 
\texttt{picture: 2, grape: 2, shoe: 0} \\
\\ 
\texttt{IMPORTANT: Your messages, unless it is the secret proposal, are directly transmitted to the other player, so do not include any response to the rules or text announcing your message. To make a secret proposal, use the indicated format. Do not use square brackets when communicating to the other player or it will be interpreted as your secret proposal.} \\
\\ 
\\ 
\texttt{The other player has submitted a secret proposal. It is now your turn to submit a proposal of your own.} \\
            }
        }
    }
    & \\ \\

    \theutterance \stepcounter{utterance}  
    & & & \multicolumn{4}{p{0.6\linewidth}}{
        \cellcolor[rgb]{0.9,0.9,0.9}{
            \makecell[{{p{\linewidth}}}]{
                \texttt{\tiny{[GM$\langle$P2]}}
                \texttt{{[}Proposal: 1 picture, 3 grapes{]}} \\
            }
        }
    }
     \\ \\

    \theutterance \stepcounter{utterance}  
    & & & \multicolumn{2}{p{0.3\linewidth}}{
        \cellcolor[rgb]{0.9,0.9,0.9}{
            \makecell[{{p{\linewidth}}}]{
                \texttt{\tiny{[GM$|$GM]}}
                \texttt{proposal} \\
            }
        }
    }
    & & \\ \\

    \theutterance \stepcounter{utterance}  
    & & & \multicolumn{2}{p{0.3\linewidth}}{
        \cellcolor[rgb]{0.9,0.9,0.9}{
            \makecell[{{p{\linewidth}}}]{
                \texttt{\tiny{[GM$|$GM]}}
                [1, 3, 0]
            }
        }
    }
    & & \\ \\

    \theutterance \stepcounter{utterance}  
    & & & \multicolumn{2}{p{0.3\linewidth}}{
        \cellcolor[rgb]{0.9,0.9,0.9}{
            \makecell[{{p{\linewidth}}}]{
                \texttt{\tiny{[GM$|$GM]}}
                [[2, 0, 1], [1, 3, 0]]
            }
        }
    }
    & & \\ \\

\end{supertabular}
}

\end{document}
